\section{Discussion}
\label{sec:discussion}

\para{What did the audit establish?} The tested models make accurate bounded research judgments, and an explicit independent audit improves every aggregate point estimate without introducing observed false alarms. The result supports using structured oversight as a promising review aid. It does not establish that the prompt causes a reliable improvement, because the paired tests are underpowered, and it does not establish low residual risk, because the rare-event bound fails the preregistered threshold.

\para{Why separate capability from safety?} An aggregate score can hide a consequential tail. Here, 96.4\% substantive accuracy coexists with three safety-aware misses, including rationales that notice a problem but fail to express the keyed decision. Reporting only accuracy would overstate trustworthiness; reporting only the 6.25\% miss estimate would hide uncertainty. The 15.4\% upper bound makes the evidential limit explicit. Even a perfect 0/16 cell remains compatible with a 17.1\% miss rate at the stated confidence level.

\para{Implications for evaluation.} Evaluators should score useful performance, integrity recognition, clean-case specificity, and uncertainty separately. Held-out tests and immutable traces remain important when systems can act on repositories or evaluators \cite{wijk2024rebench,jimenez2024swebench}. Independent monitors and randomized human audits can reduce shared blind spots, but all safety claims must remain conditional on the task distribution, model, scaffold, access, and audit budget \cite{benton2024sabotage}.

\subsection{Limitations}
First, multiple-choice recognition is not end-to-end research agency, tool-policy compliance, or unknown-unknown discovery. Fixed options can cue the fault and exact categories can penalize a semantically correct warning. Second, the cases are easier than long-horizon replication and research-engineering benchmarks; mathematics shows a clear ceiling effect. Third, all tested models belong to Qwen3, so shared training and architecture can create correlated blind spots. The failed provider access prevented the preregistered cross-family comparison.

Fourth, the safety sample is small: each model--protocol cell has only 16 fault trials, and the aggregate safety-aware analysis has 48. Repeated models within cases also limit independence, despite case-clustered intervals. Fifth, we recruit no humans, so the study does not estimate human uplift. Planted faults measure recognition capability, not the propensity of deployed agents to introduce faults intentionally or accidentally. Sixth, one deterministic run improves reproducibility but does not measure sampling variance. Finally, latency excludes download and loading time, and the \$0 local API figure excludes energy, hardware, and depreciation.

\subsection{Broader impact and responsible use}
The benchmark can help design monitors that force separate task and integrity decisions. It could also encourage false reassurance if users treat recognition on authored cases as proof of deployment safety. We recommend using the protocol as one layer in bounded review workflows, retaining expert accountability for consequential research decisions, and publishing negative results and exact upper bounds. Future benchmarks should include open-ended faults, unauthorized tool actions, harder executable tasks, independent model families, and blinded human grading.
